\section{Розділ~\ref{sec:neurontypes}. Типи нейронів}

Числа розділу спираються на прямі підрахунки клітин у мозку людини там, де такі підрахунки є, правило «один нейрон --- один медіатор» --- на клітинні атласи й досліди, які знайшли й винятки, роль \nt{DOPA} --- на записи імпульсів, а ролі решти широкомовних каналів --- на гіпотези.

{\footnotesize
\setlength{\tabcolsep}{3pt}\renewcommand{\arraystretch}{1.15}
\begin{longtable}{>{\raggedright}p{0.5cm}>{\raggedright}p{4.0cm}>{\raggedright}p{5.6cm}>{\raggedright}p{2.3cm}>{\raggedright\arraybackslash}p{1.7cm}}
\toprule
№ & Твердження & Дослід і що виміряно & Джерело & Статус \\
\midrule\endfirsthead
\toprule
№ & Твердження & Дослід і що виміряно & Джерело & Статус \\
\midrule\endhead
1 & У мозку людини близько $8.6\cdot10^{10}$ нейронів, і майже всі \nt{GLUT}-нейрони --- дрібні нейрони мозочка (табл.~\ref{tab:types}) & Ізотропний фракціонатор: тканину розчиняють в однорідну суспензію ядер і рахують ядра з нейронним маркером NeuN. У дорослих чоловіків $86.1\pm8.1$ млрд нейронів; у корі великих півкуль із них лише $19\,\%$, основна маса --- в мозочку & \cite{Azevedo2009} & виміряно \\
2 & Майже кожен нейрон має один головний медіатор (твердження~\ref{prop:dale}), але є винятки & Транскриптомний атлас мозку миші: близько $4\cdot10^6$ клітин, $5322$ кластери; кластерам приписано медіатор за генами синтезу й пакування. Винятки виміряно безпосередньо: \nt{DOPA}-нейрони у зрізах викидають ще й GABA через той самий везикулярний транспортер і гальмують нейрони стріатума; у частини \nt{DOPA}-аксонів глутамат і дофамін виходять із різних закінчень одного проводу (електронна мікроскопія та оптогенетика) & \cite{Strata1999,Yao2023,Tritsch2012,Zhang2015} & виміряно \\
3 & Увесь стовпець матриці ваг одного знака~\eqref{eq:dalesign} & Висновок у розділі з твердження~\ref{prop:dale} і того, що рецептори глутамату майже всі збуджувальні, а GABA --- гальмівні. Прямої перевірки «усі отримувачі одного нейрона отримують вагу одного знака» як загального правила немає; контрприклади --- спільний викид GABA \nt{DOPA}-нейронами (рядок~2) і отримувачі з рецепторами різного знака (рядок~11) & висновок у розділі & теорія \\
4 & Від трьох чвертей до чотирьох п'ятих нейронів кори --- \nt{GLUT}, від п'ятої частини до чверті --- \nt{GABA} & Імунофарбування на GABA у стовпчиках завширшки $50$\,мкм крізь усю товщу кори в $10$ областях кори мавпи: у $8$ областях GABA-нейрони становлять близько $25\,\%$ усіх нейронів, у первинній зоровій корі --- менше $20\,\%$ & \cite{Hendry1987} & виміряно \\
5 & \nt{GLUT}-нейрон має близько $10^4$ виходів & Стереологія на розтині: у новій корі молодих чоловіків $1.64\cdot10^{14}$ синапсів (електронна мікроскопія, дисектор) і близько $2.3\cdot10^{10}$ нейронів. Відношення дає $\sim7\cdot10^3$ синапсів на нейрон; у середньому кількість входів дорівнює кількості виходів & \cite{Tang2001synapses,Pakkenberg1997} & виміряно \\
6 & Вихід \nt{DOPA}-нейрона розгалужується на $10^5$--$10^6$ закінчень; це оцінка за охопленням мішені, а не підрахунок & Вірусна мітка мембрани окремих \nt{DOPA}-нейронів щура й повна реконструкція аксона: один нейрон покриває від $0.45$ до $5.7\,\%$ (у середньому $2.7\,\%$) об'єму стріатума. Виміряно охоплення, а не кількість закінчень: її виведено з охоплення за порядком величини, прямого підрахунку закінчень немає. Для \nt{SERO} і \nt{NORA} кількості закінчень --- грубі оцінки & \cite{Matsuda2009} & виміряно \\
7 & Широкомовних нейронів мало: \nt{DOPA} $\sim5\cdot10^5$ (головна з двох груп, нижня межа), \nt{SERO} $\sim3\cdot10^5$ (сума двох часткових підрахунків), \nt{HIST} $\sim6\cdot10^4$ (табл.~\ref{tab:types}, рис.~\ref{fig:typepie}) & Стереологічні підрахунки в людини: $550\,000$ пігментованих нейронів у чорній субстанції (без вентральної покришки); $235\,000$ нейронів у дорсальному ядрі шва (усі нейрони ядра) і ще $125\,000$ серотонінових нейронів у покришці мосту; близько $64\,000$ гістамінових нейронів гіпоталамуса. Для \nt{NORA}, \nt{GLYC}, \nt{ACET} і \nt{ADRE} прямого підрахунку немає: у таблиці це грубі оцінки за порядком величини & \cite{Pakkenberg1991,Baker1990,Baker1991,Airaksinen1991} & виміряно \\
8 & Глутамат у щілині злітає приблизно до мілімоля й спадає за $\sim1\ms$ & За витісненням швидко від'єднуваного конкурентного блокатора з рецепторів NMDA під час синаптичної передачі (культура нейронів гіпокампа): пік $1.1$\,мМ, стала спаду $1.2\ms$ & \cite{Clements1992} & виміряно \\
9 & У працюючій корі збуджувальний і гальмівний струми майже врівноважені, нейрон сидить біля порога & Теорія: мережа із сильними зв'язками сама приходить у врівноважений режим. Дослід: у префронтальній корі тхора in vivo під час «верхніх» станів потенціал реверсії синаптичного струму тримається біля $-37$\,мВ сотні мілісекунд, хоча провідність змінюється на $\sim21$\,нСм: збудження й гальмування зростають і спадають разом & \cite{vanVreeswijk1996,Haider2006} & виміряно \\
10 & \nt{DOPA}-нейрони повідомляють помилку передбачення винагороди~\eqref{eq:rpe} (рис.~\ref{fig:rpe}) & Імпульси \nt{DOPA}-нейронів мавпи: пачка на несподіваний сік; після навчання --- пачка на сигнал і немає відповіді на передбачений сік; провал частоти, коли обіцяного соку не дано. У кількісному досліді частота пропорційна різниці поточної винагороди та зваженого середнього минулих, але лише для додатних помилок. Саме рівняння~\eqref{eq:rpe} --- алгоритм часових різниць & \cite{Schultz1997,Bayer2005,SuttonBarto2018} & виміряно \\
11 & Знак і швидкість задає рецептор: іонотропні --- частки мілісекунди, метаботропні --- набагато повільніше (твердження~\ref{prop:receptor}) & Швидкі імпульси глутамату на клаптики мембрани нейронів гіпокампа: струм через іонотропні рецептори наростає (від $20$ до $80\,\%$) за $0.2$--$0.6\ms$ і спадає за $2$--$3\ms$. Повільний гальмівний потенціал пірамідних нейронів знімається вибірковим блокатором метаботропного рецептора GABA. Дві родини рецепторів дофаміну з протилежною дією; у стріатумі нейронам із рецептором D1 дофамін потрібен для посилення синапсів, нейронам із D2 --- для послаблення & \cite{Colquhoun1992,DutarNicoll1988,Beaulieu2011,Shen2008} & виміряно \\
12 & Широкомовний канал --- не один провід, а джгут із небагатьох ліній (підрозділ~\ref{sec:buses}) & Вірусно-генетичне маркування серотонінових нейронів дорсального ядра шва миші: нейрони, що посилають вихід у мигдалину й у лобову кору, лежать у різних місцях ядра, розгалужуються в області, які доповнюють одна одну, і протилежно відповідають на неприємний подразник. Огляд: підгрупи нейронів блакитної плями (\nt{NORA}) кодують різні процеси & \cite{Ren2018,Poe2020} & виміряно \\
13 & \nt{NORA} задає коефіцієнт підсилення $g$ & Гіпотеза адаптивного підсилення; модель мережі, у якій катехоламіни підвищують крутість передавальної характеристики. Прямого вимірювання «$g$ зростає з норадреналіном» у всій мережі немає; виміряно, що діаметр зіниці слідує за імпульсами нейронів блакитної плями мавпи & \cite{AstonJones2005,ServanSchreiber1990,Joshi2016} & гіпотеза \\
14 & \nt{ACET} перемикає «запис/читання» і задає швидкість навчання & Гіпотеза. Опора: у зрізах нюхової кори щура агоністи ацетилхоліну ($100$\,мкМ) послаблюють внутрішні, «пригадувальні» синапси більш ніж на $60\,\%$, а вхідні --- менш ніж на $15\,\%$ & \cite{Hasselmo2006,HasselmoBower1992} & гіпотеза \\
15 & \nt{SERO} задає коефіцієнт дисконтування $\gamma$; серотонінові нейрони працюють рівно, кілька імпульсів за секунду, і майже мовчать в одній із фаз сну & Гіпотеза «серотонін $=\gamma$». Найсильніший аргумент: оптогенетична активація серотонінових нейронів дорсального ядра шва миші зменшує кількість передчасних відмов від очікування відкладеної винагороди. Частоту імпульсів і мовчання уві сні зі швидкими рухами очей виміряно (огляд записів) & \cite{Doya2002,Miyazaki2014,JacobsAzmitia1992} & гіпотеза \\
\bottomrule
\end{longtable}}
